\documentclass[12pt]{article}

% --- Packages ---
\usepackage[utf8]{inputenc}
\usepackage[T1]{fontenc}
\usepackage[french]{babel}
\usepackage{amsmath, amssymb}
\usepackage{geometry}
\usepackage{tikz}
\usepackage[most]{tcolorbox}
\usepackage{lmodern}
\usepackage{pgfplots}
\pgfplotsset{compat=1.18}
\usepackage{mdframed}
\usepackage{bm}
\usepackage{xcolor}
\usepackage{sectsty}
\usepackage[utf8]{inputenc}
\usepackage[T1]{fontenc}
\usepackage[dvipsnames]{xcolor}
\usepackage{tkz-tab}
\usepackage{multido}

% mettre dans le préambule, après \usepackage{xcolor}
\usepackage{sectsty} % gère facilement la couleur des titres
\setlength{\parindent}{0pt}

\sectionfont{\color{red}\Large\bfseries}     % pour \section*
% \subsectionfont{\color{red}\normalsize\bfseries} % pour \subsection*


% --- Page ---
\geometry{margin=2cm}

\begin{document}

% \begin{center}
%     {\LARGE \textbf{Fonctions de référence}}\\[4mm]
%     {\large Poly à compléter – Seconde}
% \end{center}

% \hrule
% \vspace{0.4cm}

% ============================
%   PARTIE 6 : fonctions affines
% ============================

\section*{\textcolor{red}{Partie 6 : Fonctions affines}}

\subsection*{\textcolor{red}{6.1 Définition d'une fonction affine}}

\vspace{5pt}

\begin{mdframed}[
    linewidth=1pt,
    roundcorner=5pt,
    linecolor=ForestGreen
]
{\color{ForestGreen}\textbf{Définition :}}
$a$ et $b$ sont deux réels fixés. Une \textbf{fonction affine}
est une fonction définie sur $\mathbb{R}$ par :
\[
f(x) = ax + b
\]
\end{mdframed}

\underline{Exemples :}
\vspace{5pt}
\begin{itemize}
    \item $f(x) = 2x - 3$ est une fonction affine avec $a = ......$ et $b = ......$.
    \vspace{10pt}
    \item $g(x) = 3(5 - 7x) - 2(5x - 1) = $
    \dotfill
    \vspace{10pt}
    \\Donc $g$ est une fonction ................. avec $a = ......$ et $b = ......$.
    \vspace{10pt}
    \item $k(x) = -9x$ est une fonction affine avec $a = ......$ et $b = ......$.
    \vspace{10pt}
    \item $i(x) = \dfrac{-7}{2}$ est une fonction affine avec $a = ......$ et $b = ......$.
\end{itemize}

\vspace{10pt}

\underline{Remarques :}
\vspace{5pt}
\begin{itemize}
    \item Toutes les fonctions affines telles que $b = 0$ sont appelées \textbf{fonctions linéaires}.
    Elles sont de la forme $f(x) = ax$.
    \item Toutes les fonctions affines telles que $a = 0$ sont appelées \textbf{fonctions constantes}.
    Elles sont de la forme $f(x) = b$.
\end{itemize}

\vspace{10pt}

\begin{mdframed}[
    linewidth=1pt,
    roundcorner=5pt,
    linecolor=Red,
    innertopmargin=10pt,
    innerbottommargin=10pt
]
{\color{red}
\underline{Propriété :} \textbf{Taux d'accroissement d'une fonction affine}
\vspace{5pt}
\\Soient $a$ et $b$ deux réels.
\\Si $f$ est une fonction affine telle que $f(x) = ax + b$, alors pour tout $x_1$, $x_2$, on peut
déterminer le coefficient directeur $a$ :
\[
a = \dfrac{f(x_1) - f(x_2)}{x_1 - x_2}
\]
C'est le \textbf{taux d'accroissement} de la fonction $f$ entre $x_1$ et $x_2$.
}
\end{mdframed}

\underline{Démonstration :}
\vspace{5pt}
\\$\dfrac{f(x_1) - f(x_2)}{x_1 - x_2} = \dfrac{(ax_1 + b) - (ax_2 + b)}{x_1 - x_2}
= \dfrac{ax_1 + b - ax_2 - b}{x_1 - x_2} = \dfrac{a(x_1 - x_2)}{x_1 - x_2} = a$

\vspace{10pt}

\begin{mdframed}[
    linewidth=1pt,
    roundcorner=5pt,
    linecolor=RoyalBlue,
    innertopmargin=10pt,
    innerbottommargin=10pt
]
{\color{RoyalBlue}\textbf{\underline{Exercice d'application n°1 :} Savoir déterminer l'expression d'une fonction
affine}
\vspace{5pt}
\begin{enumerate}
    \item Déterminer la fonction linéaire $f$ telle que : $f(-15) = 5$.
    \item Déterminer la fonction affine $g$ telle que $g(-1) = 4$ et $g(2) = -5$.
\end{enumerate}
}
\end{mdframed}

\newpage
\subsection*{\textcolor{red}{6.2 Représentation graphique d'une fonction affine}}

\vspace{5pt}

\begin{mdframed}[
    linewidth=1pt,
    roundcorner=5pt,
    linecolor=ForestGreen,
    innerbottommargin=10pt
]
{\color{ForestGreen}\textbf{Définition :}}
    \vspace{5pt}
    \\La représentation graphique d'une fonction affine $f(x) = ax + b$ est \textbf{une droite}.
    \vspace{5pt}
    \\On dit que cette droite a pour équation $\bm{y = ax + b}$.
    \vspace{5pt}
    \\$\bm{a}$ est appelé \textbf{coefficient directeur} de la droite.
    \vspace{5pt}
    \\$\bm{b}$ est appelé \textbf{l'ordonnée à l'origine} : $f(0) = b$ donc la droite passe par le point
    $(0;b)$.
\end{mdframed}

\underline{Remarques :}
\vspace{5pt}
\begin{itemize}
    \item La représentation graphique d’une \textbf{fonction linéaire} est une droite d’équation 
    $y = ax$ qui \underline{passe par l’origine du repère}.
    \vspace{5pt}
    \item La représentation graphique d’une \textbf{fonction constante} est une droite d’équation
    $y = b$ \underline{parallèle à l’axe des abscisses}.
\end{itemize}

\vspace{10pt}

\begin{mdframed}[
    linewidth=1pt,
    roundcorner=5pt,
    linecolor=RoyalBlue,
    innertopmargin=10pt,
    innerbottommargin=10pt
]
{\color{RoyalBlue}\textbf{\underline{Exercice d'application n°2 :} Savoir déterminer l'équation
d'une droite tracée (l'expression algébrique d'une fonction affine)}}
\vspace{5pt}

\begin{mdframed}[
    linewidth=0.5pt,
    roundcorner=5pt,
    linecolor=black,
    innertopmargin=10pt,
    innerbottommargin=10pt,
    backgroundcolor=black!6
]
\underline{Méthode :}
\vspace{5pt}
\\On lit sur le graphique la valeur de $b$ : \textbf{l'ordonnée à l'origine}.
\vspace{10pt}
\\
$a = \dfrac{f(x_1) - f(x_2)}{x_1 - x_2} = \dfrac{accroissement\ vertical}{accroissement\ horizontal}
 = \dfrac{\Delta y}{\Delta x}$ : c'est le taux de variation de $f$ entre $x_1$ et $x_2$, \textbf{on lit
 donc la valeur de $a$ grâce au quadrillage}.
\end{mdframed}

\begin{center}
    \begin{tabular}{cc}
    % --- GRAPHIQUE 1 ---
    \begin{tikzpicture}[scale=1]
    \begin{axis}[
        width=8cm,
        height=8cm,
        axis lines=center,
        grid=both,
        xmin=-4, xmax=4,
        ymin=-4, ymax=4,
        xtick={-4,-3,-2,-1,0,1,2,3,4},
        ytick={-4,-3,-2,-1,0,1,2,3,4},
        samples=200,
        title={fonction n°1},
    ]
    \addplot[domain=-4:4, thick, red] {3*x - 2};
    \end{axis}
    \end{tikzpicture}
    &
    % --- GRAPHIQUE 2 ---
    \begin{tikzpicture}[scale=1]
    \begin{axis}[
        width=8cm,
        height=8cm,
        axis lines=center,
        grid=both,
        xmin=-4, xmax=4,
        ymin=-4, ymax=4,
        xtick={-4,-3,-2,-1,0,1,2,3,4},
        ytick={-4,-3,-2,-1,0,1,2,3,4},
        samples=200,
        title={fonction n°2},
    ]
    \addplot[domain=-4:4, thick, red] {-2*x + 1};
    \end{axis}
    \end{tikzpicture}
    \end{tabular}

    % ----------------------- %
    \begin{tabular}{cc}
    % --- GRAPHIQUE 3 ---
    \begin{tikzpicture}[scale=1]
    \begin{axis}[
        width=8cm,
        height=8cm,
        axis lines=center,
        grid=both,
        xmin=-4, xmax=4,
        ymin=-4, ymax=4,
        xtick={-4,-3,-2,-1,0,1,2,3,4},
        ytick={-4,-3,-2,-1,0,1,2,3,4},
        samples=200,
        title={fonction n°3},
    ]
    \addplot[domain=-4:4, thick, red] {4*x};
    \end{axis}
    \end{tikzpicture}
    &
    % --- GRAPHIQUE 4 ---
    \begin{tikzpicture}[scale=1]
    \begin{axis}[
        width=8cm,
        height=8cm,
        axis lines=center,
        grid=both,
        xmin=-4, xmax=4,
        ymin=-4, ymax=4,
        xtick={-4,-3,-2,-1,0,1,2,3,4},
        ytick={-4,-3,-2,-1,0,1,2,3,4},
        samples=200,
        title={fonction n°4},
    ]
    \addplot[domain=-4:4, thick, red] {(1/2)*x - 1};
    \end{axis}
    \end{tikzpicture}
    \end{tabular}
    % --- GRAPHIQUE 5 ---
    \begin{tikzpicture}[scale=1]
    \begin{axis}[
        width=8cm,
        height=8cm,
        axis lines=center,
        grid=both,
        xmin=-4, xmax=4,
        ymin=-2, ymax=6,
        xtick={-4,-3,-2,-1,0,1,2,3,4},
        ytick={-2,-1,0,1,2,3,4,5,6},
        samples=200,
        title={fonction n°5},
    ]
    \addplot[domain=-4:4, thick, red] {(-1/4)*x + 3};
    \end{axis}
    \end{tikzpicture}
\end{center}
\end{mdframed}

\begin{mdframed}[
    linewidth=1pt,
    roundcorner=5pt,
    linecolor=RoyalBlue,
    % innertopmargin=10pt,
    % innerbottommargin=10pt
]
{\color{RoyalBlue}\textbf{\underline{Exercice d'application n°3 :} Tracer la représentation graphique
d'une fonction affine}}
% \vspace{5pt}
\\Représenter graphiquement les fonctions suivantes :
\begin{center}
    \setlength{\tabcolsep}{15pt}
    \begin{tabular}{ccc}
        $f_1(x) = 2x - 2$
        &
        $f_2(x) = -5x + 3$
        &
        $f_3(x) = 2 - \dfrac{3}{2}x$
    \end{tabular}
\end{center}

\begin{mdframed}[
    linewidth=0.5pt,
    roundcorner=5pt,
    linecolor=black,
    % innertopmargin=10pt,
    % innerbottommargin=10pt,
    backgroundcolor=black!6
]
\underline{Méthode 1 :}
\vspace{5pt}
\\On calcule l'image de deux nombres, ce qui nous donne les coordonnées de deux points, suffisants
pour tracer la droite que l'on veut représenter.
\vspace{5pt}
\\
\underline{Méthode 2 :}
\vspace{5pt}
\\On place l'ordonnée à l'origine et on place un second point en utilisant le quadrillage et le 
coefficient $a$.
\end{mdframed}

\end{mdframed}

\end{document}
